% vim: tw=78 ai spell spelllang=en

\documentclass[10pt]{llncs}
\usepackage{color}
\usepackage{url}
\usepackage{graphicx}
\usepackage{listings}
\usepackage{hyperref}
\usepackage{xspace}
\pagestyle{plain}

\lstset{basicstyle=\ttfamily\scriptsize,escapeinside={(*@}{@*)}}

%{{{ (La)TeX definitions

\newcommand{\nb}[1]{\marginpar{\scriptsize #1}}
\renewcommand\note[1]{\nb{\textcolor{blue}{#1}}}
\newcommand\todo[1]{\nb{\textcolor{red}{#1}}}
\newcommand\tool{\textsc{Symbiotic}\xspace}
\newcommand\klee{\textsc{Klee}\xspace}

%\renewcommand\note[1]{}
%\renewcommand\todo[1]{}

%}}}

\begin{document}

%{{{ Title + Author + Abstract

\title{\tool~2: More Precise Slicing\thanks{The authors are supported by the Czech Science
    Foundation grant~P202-10-1469.}}
\subtitle{(Competition Contribution)}
\author{Jiri~Slaby \and Jan~Strej\v{c}ek}
\institute{Faculty of Informatics, Masaryk University\\
  Botanick\'a 68a, 60200 Brno, Czech Republic\\
  \email{\{slaby,strejcek\}@fi.muni.cz}}

\maketitle

\begin{abstract}
  \tool~2 keeps the concept and the structure of the original bug-finding
  tool \tool, but it uses a more precise slicing based on a field-sensitive
  pointer analysis instead of field-insensitive analysis of the original
  tool. The paper discusses this improvement and its consequences.  We also
  briefly recall basic principles of the tool, its strong and weak points,
  installation, and running instructions. Finally, we comment the results
  achieved by \tool 2 in the competition.
\end{abstract}

%}}}

%{{{ Verification approach

\section{Verification Approach and Software Architecture}
\label{sec:approach}

Both \tool~\cite{SST13} and \tool 2 implement our verification approach
proposed earlier~\cite{sse}. The approach combines three standard techniques,
namely \emph{code instrumentation}, \emph{program slicing}~\cite{weiser}, and
\emph{symbolic execution}~\cite{king}. While the approach was originally
designed for the detection of bugs described by state-machines, \tool 2 still
supports only one kind of bugs: reachability of an \texttt{ERROR} label.
Hence, we briefly recall the approach restricted just to this simple kind of
errors. We explain the structure of the tool simultaneously.
\begin{enumerate}
\item Code instrumentation inserts \texttt{assert(0)} to each \texttt{ERROR}
  label. It is performed by a \texttt{bash} script calling \texttt{sed}. The
  instrumented code is translated into the \textsc{Llvm} bitcode by
  the \textsc{Clang} compiler.
\item Program slicing removes instructions of the instrumented code that do
  not affect reachability of the inserted \texttt{assert(0)}
  statements. This code size reduction is crucial for the overall efficiency
  of \tool 2. The slicer is implemented in C++ as a plug-in for the
  \textsc{Llvm} optimizer \texttt{opt}.
\item Symbolic execution either reaches \texttt{assert(0)}, or correctly
  finishes the execution without reaching \texttt{assert(0)}, or it runs out
  of time or memory etc. These possibilities correspond to answers
  \texttt{FALSE}, \texttt{TRUE}, \texttt{UNKNOWN}, respectively.
  We use the symbolic executor \klee~\cite{klee} whose outputs are translated
  to \texttt{TRUE/FALSE/UNKNOWN} by a simple \texttt{bash} script.
\end{enumerate}
The whole pipeline is executed stepwise by another \texttt{bash} script.

All improvements of \tool 2 over the tool \tool competing in SV-COMP 2013
are in the slicer.  We have fixed some bugs in the original slicer (invalid
treatment of several instructions and functions with variable number of
arguments). While the original fixpoint algorithm for slicing~\cite{weiser}
is relatively simple, it gets more complicated for programs with pointers as
one instruction can influence the following one without any syntactic
overlap. We need to use a pointer analysis (also called points-to analysis)
to know which pointers can point to the same target. Both \tool and \tool 2
use Andersen's pointer analysis~\cite{And94}, \tool 2 replaces the original
field-insensitive analysis by a field-sensitive one. This means that every
field of a \texttt{struct} is now handled as a different pointer
target. Similarly, we handle the first 64 elements of each array as distinct
targets. The field-sensitive analysis is computationally more
demanding. Thus we have also added some type filters that speed up the
pointer analysis and make its results more accurate. All these improvements
of the slicer significantly reduce the number of incorrect answers produced
by \tool 2. For example, \tool 2 produces only correct answers for the
category \emph{ProductLines} while \tool produces 131 incorrect answers for
the same category.

%}}} 
%{{{ Discussion

\section{Strengths and Weaknesses}

Our tool is applicable to all competition benchmarks satisfying two
restrictions: the studied property is \texttt{ERROR} label reachability and
the benchmark code is a sequential C program. Hence, the results of \tool 2
in the competition categories \emph{MemorySafety} and \emph{Concurrency}
should be ignored (we missed the opt-out deadline). The first restriction
can be removed by the implementation of a more sophisticated code
instrumentation.  The second restriction comes directly from the approach as
symbolic execution and program slicing are primarily designed for sequential
programs.

\medskip We first discuss strong and weak aspects of the approach and then
we talk about additional strong and weak aspects of the tool. Our approach
is based on symbolic execution which should produce only correct answers. On
the other hand, symbolic execution suffers from the path explosion problem
and relies on expensive (and often even undecidable) SMT solving. Hence,
application of symbolic execution leads to many \texttt{UNKNOWN} answers,
which is also the main weakness of the approach. To reduce this weakness, we
combine symbolic execution with slicing which is the only theoretical source
of incorrect answers (namely false positives) of our approach. Indeed,
slicing can in some cases remove an infinite loop and a potential
unreachable \texttt{ERROR} label located below the cycle thus become
reachable. However, this situation is very rare in practice (e.g.~it does
not appear in the competition benchmarks) and we do not see it as a problem.
An orthogonal method to reduce the high cost of symbolic execution is to use
some of its variants suppressing the path explosion problem. For example, we
plan to apply \emph{compact symbolic execution}~\cite{CSE} instead of the
classic one.

The strong aspect of the tool is its simple architecture: it is a sequence
of scripts and standalone tools that are easy to replace (for example, if
there is a better symbolic executor for \textsc{Llvm} bitcode, we can deploy
it in few minutes). The main weakness of \tool 2 lies in the incorrect
results which are sometimes reported.  Even if the number of incorrect
results is substantially lower than in the case of \tool, it is still
relatively high.  All the incorrect results are due to imperfection of our
implementation.

%}}}
%{{{ Tool setup and configuration

\section{Tool Setup and Configuration}
Before using \tool 2, ensure that the target system contains 32-bit
libraries (for 32-bit benchmarks) and \textsc{Llvm} with \textsc{Clang}.
\textsc{Llvm} and \textsc{Clang} have to be in version 3.2 exactly.  Then,
\tool 2 can be downloaded from \url{http://sf.net/projects/symbiotic/}.
Due to a bug in \klee causing absolute paths to be built in, \klee requires to
be run from a pre-defined path. Hence we are obliged to change the current
directory to \texttt{/opt/} and \emph{untar} the downloaded \tool 2 archive
there.  Running the tool is then straightforward.  When the current directory
is \texttt{/opt/symbiotic/}, the tool can be invoked for each
\texttt{<benchmark.c>} from the set by \texttt{./runme <benchmark.c>}. For
benchmarks intended for 64-bit, set \texttt{MFLAG=-m64} environment variable.
The answers provided by \tool 2 are as required by the competition rules:
\texttt{TRUE/FALSE/UNKWNOWN}. If the result for \texttt{<benchmark.c>} is
\texttt{FALSE}, discovered error paths can be found in
\texttt{<benchmark.c>-klee-out/}.

%}}}
%{{{ Software project and contributors

\section{Software Project and Contributors}
\tool 2 was contributed mostly by the authors of this paper and Marek
Trt\'{\i}k. Jiri Slaby is a contact person. The tool is licensed under the GNU
GPLv2 License unless specified otherwise for its parts.

% We would like to thank our advisors prof. Anton\'in Ku\v{c}era and dr. Jan
% Strej\v{c}ek. Our work is supported by grant no.

%}}}

\bibliographystyle{plain}
\bibliography{sse2014}

\end{document}
